\subsection{The literal Liu--Rollenske scalar and its double cover}
\begin{statement}{P3K.36}{Proved}
In the supplied Liu--Rollenske construction, retain the exact scalar
\(u\) in its line action \(B_u=uB_0\), its good fourth root at the
\(III^*\) source fibre, and its specified logarithmic vector
\(a=e_4/4\). On the universal logarithmic cover of a sufficiently small
punctured \(u\)-disc, the construction gives a holomorphic family of
compact complex threefolds. Its analytic continuation through one
positive \(u\)-turn is \textbf{not biholomorphic to its initial fibre}.
Consequently this fixed-data construction does not descend to the
punctured \(u\)-disc with these fixed data. The precise error is the treatment of
\(e_4/4\) as a single-valued torsion section when \(e_4\) is the
logarithm of the scalar multiplier.

The degree-two cover
\[
                         u=w^2                                      \tag{LR.1}
\]
does carry the actual holomorphic family, with the exact endpoint map
of K35. In its simultaneously closed Cox marking the positive
\(w\)-loop has
\[
 (k_{\rm Jin},\ell)=(-3,-12),\qquad
 (m_{\rm rel},n_{\rm rel},A)=(3,0,0).                 \tag{LR.2}
\]
For every allowed index-one fan, with physical delta norm
\(S=p^2-pq+q^2\), its class is
\[
       [\Sigma_{\Phi_w,T}]
              =\frac38(S-1)(S-3)\pmod {28}.           \tag{LR.3}
\]
The source's square-diagonal fan retains the physical vertical edge,
so \(S=1\) and this corrected family's return is zero. No return is
assigned to the nonclosed one-\(u\) path.

The same local obstruction holds with every retained constant root
line character: that character changes a fixed sign below, not the
branch change under one \(u\)-turn.

\end{statement}
\begin{proof}
\par

\medskip\noindent\textbf{1. The actual central torus and its scalar}

\par\nobreak

The source's equation (2.3) is the lift of translation by
\(-2P\), multiplied by \(u\). On the minimal good root the marked
section specializes to \((1+i)/2\) on
\(E=\mathbb C/(\mathbb Z+i\mathbb Z)\). Its line is the nonzero
order-two line with character
\[
                  m+in\longmapsto(-1)^{m+n}.          \tag{LR.4}
\]
These are the explicit statements and coordinate maps in Lemma 2.7 and Corollary 2.8; the same marked root
and corrected line are independently constructed in F41/F47.

Since \(2P\) specializes to zero, the restriction of \(B_0\) to
this line complement is multiplication by a constant \(c\ne0\).
The constant is independent of \(u\): every source, line, and root
comparison was fixed before multiplying the action by \(u\).
Thus its actual scalar multiplier is
\[
              \alpha=cu,\qquad
              \tau=\frac{\log(cu)}{2\pi i},
              \qquad\operatorname{Im}\tau>0.         \tag{LR.5}
\]
The line can be presented on the universal cover of \(E\) as
\((z,v)\sim(z+m+in,(-1)^{m+n}v)\). Writing
\(v=e^{2\pi i\eta}\), and then quotienting by \(B_u\), gives
exactly the central abelian surface
\[
 A_\tau=\mathbb C^2/\Lambda_\tau,
 \quad
 \Lambda_\tau=
 \left\langle
  (2,0),(1+i,0),(1,1/2),(0,\tau)
 \right\rangle_{\mathbb Z}.                           \tag{LR.6}
\]
The displayed lattice map is literal: its generators are the two
periods in the kernel of (LR.4), the remaining character period,
and the actual scalar quotient period. In particular
\[
 \delta=(0,1)=2(1,1/2)-(2,0),\qquad e_4=(0,\tau).
\]

The root deck acts on the source by \(z\mapsto-iz\). In this flat
line presentation an actual lift is
\((z,\eta)\mapsto(-iz,\eta+\kappa)\), with
\(4\kappa\in\mathbb Z\). Indeed (LR.4) is invariant under
\(-i\); the ratio of any two lifts is a nowhere-zero function on
the compact elliptic curve, hence a scalar; its fourth power is one.
The normalized source lift has \(\kappa=0\). Keeping \(\kappa\)
records every fixed finite line character. The logarithmic vector
is exactly \(a=(0,\tau/4)\), by the source's Definition 3.7. Therefore the reduced multiple fibre is
\[
 B_{\tau,\kappa}=A_\tau/\langle R\rangle,
 \qquad R(z,\eta)=(-iz,\eta+\tau/4+\kappa).           \tag{LR.7}
\]
This action has order four and is free: a nontrivial power has a
nonintegral multiple of \(\operatorname{Im}\tau\) in its second
coordinate, whereas the imaginary second-coordinate periods of
\(\Lambda_\tau\) are integral multiples of
\(\operatorname{Im}\tau\). Its fourth power is the lattice period
\((0,\tau+4\kappa)\).

For reference, the linear deck in the ordered lattice basis of
(LR.6) is
\[
 D=\begin{pmatrix}
 1&1&0&0\\-2&-1&-1&0\\0&0&1&0\\0&0&0&1
 \end{pmatrix}.
\]
Its invariant lattice is exactly
\(\mathbb Z\delta\oplus\mathbb Z e_4\), agreeing with the
source's explicitly calculated Proposition 3.6.
A positive \(u\)-turn has the exact comparison
\[
 \tau\longmapsto\tau+1,
 \quad \Lambda_{\tau+1}=\Lambda_\tau,
 \quad a\longmapsto a+\delta/4.                      \tag{LR.8}
\]
The unchanged abelian surface in (LR.8) does not identify the affine
quotients: the torsion translation in (LR.7) has changed.

\par

\medskip\noindent\textbf{2. An intrinsic invariant of the reduced fibre}

\par\nobreak

The actual holomorphic projection
\[
 \pi_{\tau,\kappa}:B_{\tau,\kappa}\longrightarrow
 E_{\tau,\kappa}:=
 \frac{\mathbb C}{\frac12\mathbb Z+
                       (\frac\tau4+\kappa)\mathbb Z},
 \qquad [(z,\eta)]\longmapsto[\eta]                  \tag{LR.9}
\]
is well-defined by every generator in (LR.6) and (LR.7).
It has connected fibres. More explicitly, the image of the full
universal-cover deck group under its second-coordinate translation
is precisely the displayed rank-two lattice; after choosing a lift
of a target point, its zero-second-coordinate subgroup has elliptic
kernel generated by \((2,0),(1+i,0)\). This is the actual fibre of
(LR.9), so it is connected.

We also verify that (LR.9) is the Albanese map, without a choice of
marking. Holomorphic one-forms on \(A_\tau\) are constant linear
forms: the coefficients of their pullback to \(\mathbb C^2\) are
periodic entire functions and hence constant by compactness. The
\(R\)-invariant subspace is exactly \(\mathbb C\,d\eta\), since
\(R^*dz=-i\,dz\) and \(R^*d\eta=d\eta\). The periods of
\(d\eta\) on the full fundamental group of \(B_{\tau,\kappa}\)
are exactly \(\frac12\mathbb Z+(\frac\tau4+\kappa)\mathbb Z\):
this follows directly from the same affine deck generators. The
integration definition of the Albanese torus therefore gives
precisely (LR.9), up to a target translation. In particular its
elliptic-curve isomorphism class is intrinsic to \(B_{\tau,\kappa}\).

Normalize its first period to one by \(\eta\mapsto2\eta\). Its
Tate parameter is
\[
 Q=\exp\bigl(2\pi i(\tau/2+2\kappa)\bigr)
    =e^{4\pi i\kappa}\sqrt{cu}.                     \tag{LR.10}
\]
The factor \(e^{4\pi i\kappa}\) is a fixed sign. Under (LR.8)
this parameter changes to \(-Q\).

For completeness, the classical Weierstrass lattice invariant has
\(j(Q)=Q^{-1}+744+O(Q)\). The leading term alone suffices and can
be obtained directly from the convergent lattice sums: summing the
geometric series in the cotangent partial fraction gives
\[
 E_4(Q)=1+240\sum_{n\ge1}\sigma_3(n)Q^n,
 \quad E_6(Q)=1-504\sum_{n\ge1}\sigma_5(n)Q^n,
 \quad
 j(Q)=\frac{1728E_4(Q)^3}{E_4(Q)^3-E_6(Q)^2}.
\]
Their denominator is \(1728Q+O(Q^2)\), so \(Qj(Q)\) extends
holomorphically with value one at zero. Consequently
\[
 Q\bigl(j(Q)-j(-Q)\bigr)
\]
extends holomorphically with value two. It is nonzero on a
sufficiently small disc. Thus
\[
            B_{\tau,\kappa}\not\simeq
            B_{\tau+1,\kappa}                       \tag{LR.11}
\]
for every sufficiently small nonzero \(u\), even without preserving
any of their markings. A biholomorphism would induce an isomorphism
of their Albanese elliptic curves, whereas (LR.10) has different
Weierstrass invariant after the turn.

\par

\medskip\noindent\textbf{3. Every threefold endpoint must preserve this fibre}

\par\nobreak

We give the intrinsic argument independently of any canonical-bundle
formula. The actual source/cap Mayer--Vietoris calculation F41/F47
proves \(H^2(Y;\mathbb R)=0\) for these sphere members. Let
\(f:Y\to\mathbb P^1\) be their constructed connected torus fibration.
Every effective divisor on \(Y\) is vertical. Otherwise a horizontal
irreducible component meets a general smooth torus fibre \(F\) in
a nonzero effective divisor \(C\). Its real first Chern class is
zero, because it is the restriction of the corresponding class on
\(Y\). Every compact complex torus admits a translation-invariant
Kähler form \(\omega\), and the divisor identity gives
\(\int_F c_1(\mathcal O(C))\wedge\omega=\int_C\omega>0\), a
contradiction.

Let \(g\) be a meromorphic function on \(Y\). Its poles are
therefore vertical. On a general fibre it has no divisorial poles;
normality and Hartogs extension remove the remaining codimension-two
indeterminacy. It is a holomorphic function on a compact connected
torus and is constant. This gives a meromorphic function on the
regular part of the base. Near a bad base value \(p\), multiply
\(g\) by a sufficiently large integer power of \(f^*(t-p)\) to
cancel every polar order on the finitely many fibre components.
The product is holomorphic by normality. It is constant on compact
connected fibres and hence descends holomorphically: its descended
function is bounded near \(p\) by properness, and the removable
singularity theorem applies. Dividing back proves
\[
                   \mathbb C(Y)=f^*\mathbb C(t).      \tag{LR.12}
\]

A biholomorphism between two such threefolds induces an automorphism
of this meromorphic function field, hence a Möbius map of the base.
Indeed a rational function with a rational compositional inverse
has degree one. Thus it fits into an actual commutative square with
their maps \(f\). Multiplicities of fibres are preserved, because a
Möbius map has nonzero local derivative. The unique multiple fibre
has multiplicity four, so its reduced support must map biholomorphically
to the unique reduced multiple fibre in the target. This contradicts
(LR.11) after a single positive \(u\)-turn. This argument excludes
all possible holomorphic endpoint maps, including source/base
symmetries and source inversion; no assumed translation normal form
for an endpoint is needed.

\par

\medskip\noindent\textbf{4. Exact correction and monodromy of the corrected family}

\par\nobreak

The earlier source choice of \(\tau\) with
\(\alpha=e^{2\pi i\tau}\), and then of \(e_4/4\), is valid on a
simply connected scalar chart. It cannot be continued single valued
through the punctured \(u\)-disc: the transition is (LR.8). Formula
(LR.11) shows that this is an actual failure of descent, rather than
only a missing chosen marking.

Apply K35 to the actual row
\[
 (a,b)=(1,2),\quad m_{III^*}=4,\quad m_{II}=1,
 \quad n_{III^*}=1,
\]
with its regular order-one datum retained. We choose the constant
normalization of \(B_0\) to be the source's actual \(B_0\), so
\(\lambda=u\). The canonical parameter is \(u=r^4\), and K35's
explicit quotient gives \(w=r^2\), exactly (LR.1). Its source map is
the meromorphic \(+G\) chord lift, made invertible on the actual
cap by the displayed root powers; its good-root deck commutator is
\(-1\). Its square is the actual one-gamma gluing. Thus the double
turn is a proved endpoint map, with the original line character,
root chart and cusp fan retained. The obstruction in Sections 2--3
proves degree two is minimal.

The actual degree-four page has \(\chi=1\) and line square \(-2\),
so K30 gives \((-6,-24)\) on the canonical \(r\)-loop. K35's
constructed twofold mapping-torus cover and closed Cox product give
(LR.2) by integral characteristic naturality. Substitution into K31
gives (LR.3). This uses the actual degree-two parameter quotient and
never divides a class in \(\Theta_7\).

The code \path{code/part3/liu_rollenske_parameter.py} checks the exact
lattice deck, invariant subspace, parameter shear, affine order-four
norm, Albanese period generators, convergent-series leading invariant,
and the corrected integer/fan identities. Its JSON and captured output
are saved. The compactness, gluing, and smoothness of the corrected
family are the actual K35 descent of F47, with all hypotheses identified
above.

\end{proof}
